\documentclass[11pt,a4paper]{article}
\usepackage[utf8]{inputenc}
\usepackage[T1]{fontenc}
\usepackage[catalan]{babel}
\usepackage{lmodern}
\usepackage[sfdefault,lining]{FiraSans}
\usepackage[a4paper,top=2.3cm,bottom=2.2cm,left=1.8cm,right=1.8cm,headheight=1.3cm,headsep=0.4cm,footskip=1cm]{geometry}
\usepackage{amsmath,amssymb}
\usepackage{icomma}
\usepackage{graphicx}
\usepackage[table]{xcolor}
\usepackage{tikz}
\usepackage{fontawesome5}
\usepackage{ccicons}
\usepackage{fancyhdr}
\usepackage{lastpage}
\usepackage{tasks}
\usepackage{needspace}
\usepackage[most]{tcolorbox}
\usepackage{array}
\usepackage[hidelinks,pdftitle={Fitxa d'exercicis de la unitat 1: Nombres naturals i enters. 2n d'ESO},pdfsubject={Matemàtiques 2n ESO},pdfkeywords={nombres naturals, divisió entera, múltiples, divisors, nombres primers, mcm, MCD, nombres enters, operacions combinades, potències, arrels}]{hyperref}

\definecolor{navy}{HTML}{1B2748}
\definecolor{accent}{HTML}{E8603A}
\definecolor{teal}{HTML}{1F9E90}
\definecolor{soft}{HTML}{EEF1F8}
\definecolor{muted}{HTML}{5D6685}

\setlength{\parindent}{0pt}
\setlength{\parskip}{3pt}
\settasks{label=\alph*), label-format=\bfseries\color{navy}, label-width=1.3em, item-indent=1.8em, column-sep=0.6em, after-item-skip=2pt, before-skip=3pt, after-skip=4pt}

% ---------- capçalera i peu ----------
\newcommand{\llicencia}{\href{https://creativecommons.org/licenses/by-nc-sa/4.0/deed.ca}{\ccbyncsa}\ \ Llicència Creative Commons Reconeixement-NoComercial-CompartirIgual 4.0 Internacional (CC BY-NC-SA 4.0)}
\pagestyle{fancy}
\fancyhf{}
\fancyhead[L]{\raisebox{-0.35\height}{\includegraphics[height=0.95cm]{logo-gva.png}}}
\fancyhead[C]{\small\color{muted}Fitxa d'exercicis · Unitat 1: Nombres naturals i enters · 2n d'ESO}
\fancyhead[R]{\raisebox{-0.35\height}{\includegraphics[height=0.95cm]{logo-vermellar.png}}}
\fancyfoot[L]{\scriptsize\color{muted}\llicencia}
\fancyfoot[R]{\small\color{muted}\thepage}
\renewcommand{\headrulewidth}{0.4pt}
\renewcommand{\headrule}{\color{muted!40}\hrule height \headrulewidth}
\fancypagestyle{primera}{\fancyhf{}\renewcommand{\headrulewidth}{0pt}
  \fancyfoot[L]{\scriptsize\color{muted}\llicencia}
  \fancyfoot[R]{\small\color{muted}\thepage}}

% ---------- nivells i exercicis ----------
\newcommand{\estrelles}[1]{{\color{accent}\ifcase#1\or\faStar\faStar[regular]\faStar[regular]\or\faStar\faStar\faStar[regular]\or\faStar\faStar\faStar\fi}}
\newcounter{exer}
\newcommand{\ex}[1]{\par\needspace{4\baselineskip}\medskip\stepcounter{exer}\noindent{\color{navy}\bfseries\theexer.}\ \ {\small\estrelles{#1}}\ \ }

\newcommand{\seccio}[2]{\par\needspace{10\baselineskip}\bigskip
  \begin{tcolorbox}[colback=navy,colframe=navy,coltext=white,arc=2mm,boxrule=0pt,left=3mm,right=3mm,top=1.2mm,bottom=1.2mm]
  {\large\bfseries #1\quad #2}
  \end{tcolorbox}\vspace{-2pt}}

\newcolumntype{C}[1]{>{\centering\arraybackslash}p{#1}}

% quadre de recordatori (només a la fitxa d'exercicis)
\newcommand{\recorda}[1]{\par\smallskip\begin{tcolorbox}[colback=accent!7,colframe=accent!55,arc=2mm,boxrule=0.5pt,left=2.5mm,right=2.5mm,top=1mm,bottom=1mm]\small{\color{accent}\faLightbulb[regular]}\ \textbf{Recorda:}\ #1\end{tcolorbox}}
% recta numèrica entera: \rectabuida{min}{max} i \recta{min}{max}{punts}
\newcommand{\rectabase}[2]{\draw[latex-latex,muted] (#1-0.8,0)--(#2+0.8,0);
  \foreach \i in {#1,...,#2}{\draw[muted] (\i,0.12)--(\i,-0.12) node[below=-1pt,font=\scriptsize,text=muted]{$\i$};}}
\newcommand{\rectabuida}[2]{\begin{tikzpicture}[x=0.78cm,baseline=0]\rectabase{#1}{#2}\end{tikzpicture}}
\newcommand{\recta}[3]{\begin{tikzpicture}[x=0.78cm,baseline=0]\rectabase{#1}{#2}
  \foreach \p in {#3}{\fill[accent] (\p,0) circle (0.11) node[above=1pt,font=\small\bfseries,text=navy]{$\p$};}\end{tikzpicture}}
\newcommand{\res}[1]{#1}

% =====================================================================
% Capçalera FentMates de les fitxes: un color per a cada curs.
% Els colors són els de la pàgina inicial (LEVELS a assets/js/materials.js);
% el logo ve de _eines/logo/pdf/ (python3 _eines/logo/genera.py).
%
% Ús, al final del preàmbul de la unitat:
%   % =====================================================================
% Capçalera FentMates de les fitxes: un color per a cada curs.
% Els colors són els de la pàgina inicial (LEVELS a assets/js/materials.js);
% el logo ve de _eines/logo/pdf/ (python3 _eines/logo/genera.py).
%
% Ús, al final del preàmbul de la unitat:
%   % =====================================================================
% Capçalera FentMates de les fitxes: un color per a cada curs.
% Els colors són els de la pàgina inicial (LEVELS a assets/js/materials.js);
% el logo ve de _eines/logo/pdf/ (python3 _eines/logo/genera.py).
%
% Ús, al final del preàmbul de la unitat:
%   \input{../comu/fentmates.tex}
%   \fmcurs{1eso}
% i, al començament del document:
%   \capcalera{Matemàtiques}{1r ESO}{Fitxa d'exercicis · Unitat 2}{Divisibilitat}
% Opcional abans de \input: \def\fmneed{14} (espai mínim abans de cada secció).
% =====================================================================
\makeatletter
\def\fm@logodir{../../logo/pdf}
\providecommand{\fmneed}{10}
\def\fm@cursnom{}

% clau del curs, color i color del text damunt del color
\newcommand{\fm@curs}[3]{\@namedef{fm@color@#1}{#2}\@namedef{fm@text@#1}{#3}}
\fm@curs{1eso}{E6B84A}{navy}
\fm@curs{2eso}{3B82F6}{white}
\fm@curs{3eso}{FF7A50}{white}
\fm@curs{4eso-a}{A78BFA}{white}
\fm@curs{4eso-b}{6D28D9}{white}
\fm@curs{1batx-mi}{9B1C1C}{white}
\fm@curs{1batx-ccss}{EF6B6B}{white}
\fm@curs{2batx-mii}{1E6B47}{white}
\fm@curs{2batx-ccss}{2BB3A3}{white}

\newcommand{\fmcurs}[1]{%
  \@ifundefined{fm@color@#1}{\PackageError{fentmates}{Curs desconegut: #1}{}}{}%
  \edef\fm@tmp{\noexpand\definecolor{curs}{HTML}{\@nameuse{fm@color@#1}}%
               \noexpand\colorlet{curstext}{\@nameuse{fm@text@#1}}}\fm@tmp
  \colorlet{cursfosc}{curs!72!black}%
  \def\fm@logo{\fm@logodir/vertical-#1.pdf}%
  \def\fm@marca{\fm@logodir/marca-#1.pdf}}

% ---------- capçalera de la primera pàgina ----------
% #1 matèria (petit), #2 curs, #3 tipus de document i unitat, #4 títol
\newcommand{\capcalera}[4]{%
  \gdef\fm@cursnom{#2}%
  \par\noindent
  \begin{tcolorbox}[colback=curs,colframe=curs,boxrule=0pt,arc=2.2mm,width=3.5cm,height=2.15cm,
      valign=center,left=3mm,right=2mm,top=1mm,bottom=1mm,nobeforeafter,box align=center]
    \color{curstext}{\footnotesize #1}\par\vspace{1pt}{\fontsize{19}{21}\selectfont\bfseries #2\par}%
  \end{tcolorbox}\hspace{5mm}%
  \begin{minipage}[c]{\dimexpr\linewidth-3.5cm-5mm-2.9cm\relax}
    {\color{cursfosc}\small\bfseries\MakeUppercase{#3}}\par\vspace{2pt}
    {\color{navy}\fontsize{20}{23}\selectfont\bfseries #4\par}
  \end{minipage}\hfill
  \begin{minipage}[c]{2.6cm}\raggedleft\includegraphics[width=2.6cm]{\fm@logo}\end{minipage}\par
  \vspace{2pt}}

% ---------- barres de secció del color del curs ----------
\renewcommand{\seccio}[2]{\par\needspace{\fmneed\baselineskip}\bigskip
  \begin{tcolorbox}[colback=curs,colframe=curs,coltext=curstext,arc=2mm,boxrule=0pt,left=3mm,right=3mm,top=1.2mm,bottom=1.2mm]
  {\large\bfseries #1\quad #2}
  \end{tcolorbox}\vspace{-2pt}}

% ---------- estrelles de dificultat del color del curs (a Batxillerat no se'n posen) ----------
\@ifundefined{estrelles}{}{%
  \renewcommand{\estrelles}[1]{{\color{curs}\ifcase#1\or\faStar\faStar[regular]\faStar[regular]\or\faStar\faStar\faStar[regular]\or\faStar\faStar\faStar\fi}}}

% ---------- capçalera i peu de les pàgines ----------
% primera pàgina: logos institucionals al peu
\fancypagestyle{primera}{\fancyhf{}\renewcommand{\headrulewidth}{0pt}
  \fancyfoot[L]{\raisebox{-0.32\height}{\includegraphics[height=0.78cm]{logo-gva.png}}\hspace{5mm}%
    \raisebox{-0.32\height}{\includegraphics[height=0.78cm]{logo-vermellar.png}}}
  \fancyfoot[R]{\scriptsize\color{muted}\href{https://creativecommons.org/licenses/by-nc-sa/4.0/deed.ca}{\ccbyncsa}\ \ CC BY-NC-SA 4.0 · fentmates.org\qquad\small\thepage}}
% la resta: el símbol a l'esquerra i el curs a la dreta, amb el títol al centre com abans
\fancyhead[L]{\raisebox{-0.3\height}{\includegraphics[height=0.62cm]{\fm@marca}}}
\fancyhead[R]{\small\bfseries\color{cursfosc}\fm@cursnom}
\renewcommand{\headrule}{\color{curs!45}\hrule height \headrulewidth}
\makeatother

%   \fmcurs{1eso}
% i, al començament del document:
%   \capcalera{Matemàtiques}{1r ESO}{Fitxa d'exercicis · Unitat 2}{Divisibilitat}
% Opcional abans de \input: \def\fmneed{14} (espai mínim abans de cada secció).
% =====================================================================
\makeatletter
\def\fm@logodir{../../logo/pdf}
\providecommand{\fmneed}{10}
\def\fm@cursnom{}

% clau del curs, color i color del text damunt del color
\newcommand{\fm@curs}[3]{\@namedef{fm@color@#1}{#2}\@namedef{fm@text@#1}{#3}}
\fm@curs{1eso}{E6B84A}{navy}
\fm@curs{2eso}{3B82F6}{white}
\fm@curs{3eso}{FF7A50}{white}
\fm@curs{4eso-a}{A78BFA}{white}
\fm@curs{4eso-b}{6D28D9}{white}
\fm@curs{1batx-mi}{9B1C1C}{white}
\fm@curs{1batx-ccss}{EF6B6B}{white}
\fm@curs{2batx-mii}{1E6B47}{white}
\fm@curs{2batx-ccss}{2BB3A3}{white}

\newcommand{\fmcurs}[1]{%
  \@ifundefined{fm@color@#1}{\PackageError{fentmates}{Curs desconegut: #1}{}}{}%
  \edef\fm@tmp{\noexpand\definecolor{curs}{HTML}{\@nameuse{fm@color@#1}}%
               \noexpand\colorlet{curstext}{\@nameuse{fm@text@#1}}}\fm@tmp
  \colorlet{cursfosc}{curs!72!black}%
  \def\fm@logo{\fm@logodir/vertical-#1.pdf}%
  \def\fm@marca{\fm@logodir/marca-#1.pdf}}

% ---------- capçalera de la primera pàgina ----------
% #1 matèria (petit), #2 curs, #3 tipus de document i unitat, #4 títol
\newcommand{\capcalera}[4]{%
  \gdef\fm@cursnom{#2}%
  \par\noindent
  \begin{tcolorbox}[colback=curs,colframe=curs,boxrule=0pt,arc=2.2mm,width=3.5cm,height=2.15cm,
      valign=center,left=3mm,right=2mm,top=1mm,bottom=1mm,nobeforeafter,box align=center]
    \color{curstext}{\footnotesize #1}\par\vspace{1pt}{\fontsize{19}{21}\selectfont\bfseries #2\par}%
  \end{tcolorbox}\hspace{5mm}%
  \begin{minipage}[c]{\dimexpr\linewidth-3.5cm-5mm-2.9cm\relax}
    {\color{cursfosc}\small\bfseries\MakeUppercase{#3}}\par\vspace{2pt}
    {\color{navy}\fontsize{20}{23}\selectfont\bfseries #4\par}
  \end{minipage}\hfill
  \begin{minipage}[c]{2.6cm}\raggedleft\includegraphics[width=2.6cm]{\fm@logo}\end{minipage}\par
  \vspace{2pt}}

% ---------- barres de secció del color del curs ----------
\renewcommand{\seccio}[2]{\par\needspace{\fmneed\baselineskip}\bigskip
  \begin{tcolorbox}[colback=curs,colframe=curs,coltext=curstext,arc=2mm,boxrule=0pt,left=3mm,right=3mm,top=1.2mm,bottom=1.2mm]
  {\large\bfseries #1\quad #2}
  \end{tcolorbox}\vspace{-2pt}}

% ---------- estrelles de dificultat del color del curs (a Batxillerat no se'n posen) ----------
\@ifundefined{estrelles}{}{%
  \renewcommand{\estrelles}[1]{{\color{curs}\ifcase#1\or\faStar\faStar[regular]\faStar[regular]\or\faStar\faStar\faStar[regular]\or\faStar\faStar\faStar\fi}}}

% ---------- capçalera i peu de les pàgines ----------
% primera pàgina: logos institucionals al peu
\fancypagestyle{primera}{\fancyhf{}\renewcommand{\headrulewidth}{0pt}
  \fancyfoot[L]{\raisebox{-0.32\height}{\includegraphics[height=0.78cm]{logo-gva.png}}\hspace{5mm}%
    \raisebox{-0.32\height}{\includegraphics[height=0.78cm]{logo-vermellar.png}}}
  \fancyfoot[R]{\scriptsize\color{muted}\href{https://creativecommons.org/licenses/by-nc-sa/4.0/deed.ca}{\ccbyncsa}\ \ CC BY-NC-SA 4.0 · fentmates.org\qquad\small\thepage}}
% la resta: el símbol a l'esquerra i el curs a la dreta, amb el títol al centre com abans
\fancyhead[L]{\raisebox{-0.3\height}{\includegraphics[height=0.62cm]{\fm@marca}}}
\fancyhead[R]{\small\bfseries\color{cursfosc}\fm@cursnom}
\renewcommand{\headrule}{\color{curs!45}\hrule height \headrulewidth}
\makeatother

%   \fmcurs{1eso}
% i, al començament del document:
%   \capcalera{Matemàtiques}{1r ESO}{Fitxa d'exercicis · Unitat 2}{Divisibilitat}
% Opcional abans de \input: \def\fmneed{14} (espai mínim abans de cada secció).
% =====================================================================
\makeatletter
\def\fm@logodir{../../logo/pdf}
\providecommand{\fmneed}{10}
\def\fm@cursnom{}

% clau del curs, color i color del text damunt del color
\newcommand{\fm@curs}[3]{\@namedef{fm@color@#1}{#2}\@namedef{fm@text@#1}{#3}}
\fm@curs{1eso}{E6B84A}{navy}
\fm@curs{2eso}{3B82F6}{white}
\fm@curs{3eso}{FF7A50}{white}
\fm@curs{4eso-a}{A78BFA}{white}
\fm@curs{4eso-b}{6D28D9}{white}
\fm@curs{1batx-mi}{9B1C1C}{white}
\fm@curs{1batx-ccss}{EF6B6B}{white}
\fm@curs{2batx-mii}{1E6B47}{white}
\fm@curs{2batx-ccss}{2BB3A3}{white}

\newcommand{\fmcurs}[1]{%
  \@ifundefined{fm@color@#1}{\PackageError{fentmates}{Curs desconegut: #1}{}}{}%
  \edef\fm@tmp{\noexpand\definecolor{curs}{HTML}{\@nameuse{fm@color@#1}}%
               \noexpand\colorlet{curstext}{\@nameuse{fm@text@#1}}}\fm@tmp
  \colorlet{cursfosc}{curs!72!black}%
  \def\fm@logo{\fm@logodir/vertical-#1.pdf}%
  \def\fm@marca{\fm@logodir/marca-#1.pdf}}

% ---------- capçalera de la primera pàgina ----------
% #1 matèria (petit), #2 curs, #3 tipus de document i unitat, #4 títol
\newcommand{\capcalera}[4]{%
  \gdef\fm@cursnom{#2}%
  \par\noindent
  \begin{tcolorbox}[colback=curs,colframe=curs,boxrule=0pt,arc=2.2mm,width=3.5cm,height=2.15cm,
      valign=center,left=3mm,right=2mm,top=1mm,bottom=1mm,nobeforeafter,box align=center]
    \color{curstext}{\footnotesize #1}\par\vspace{1pt}{\fontsize{19}{21}\selectfont\bfseries #2\par}%
  \end{tcolorbox}\hspace{5mm}%
  \begin{minipage}[c]{\dimexpr\linewidth-3.5cm-5mm-2.9cm\relax}
    {\color{cursfosc}\small\bfseries\MakeUppercase{#3}}\par\vspace{2pt}
    {\color{navy}\fontsize{20}{23}\selectfont\bfseries #4\par}
  \end{minipage}\hfill
  \begin{minipage}[c]{2.6cm}\raggedleft\includegraphics[width=2.6cm]{\fm@logo}\end{minipage}\par
  \vspace{2pt}}

% ---------- barres de secció del color del curs ----------
\renewcommand{\seccio}[2]{\par\needspace{\fmneed\baselineskip}\bigskip
  \begin{tcolorbox}[colback=curs,colframe=curs,coltext=curstext,arc=2mm,boxrule=0pt,left=3mm,right=3mm,top=1.2mm,bottom=1.2mm]
  {\large\bfseries #1\quad #2}
  \end{tcolorbox}\vspace{-2pt}}

% ---------- estrelles de dificultat del color del curs (a Batxillerat no se'n posen) ----------
\@ifundefined{estrelles}{}{%
  \renewcommand{\estrelles}[1]{{\color{curs}\ifcase#1\or\faStar\faStar[regular]\faStar[regular]\or\faStar\faStar\faStar[regular]\or\faStar\faStar\faStar\fi}}}

% ---------- capçalera i peu de les pàgines ----------
% primera pàgina: logos institucionals al peu
\fancypagestyle{primera}{\fancyhf{}\renewcommand{\headrulewidth}{0pt}
  \fancyfoot[L]{\raisebox{-0.32\height}{\includegraphics[height=0.78cm]{logo-gva.png}}\hspace{5mm}%
    \raisebox{-0.32\height}{\includegraphics[height=0.78cm]{logo-vermellar.png}}}
  \fancyfoot[R]{\scriptsize\color{muted}\href{https://creativecommons.org/licenses/by-nc-sa/4.0/deed.ca}{\ccbyncsa}\ \ CC BY-NC-SA 4.0 · fentmates.org\qquad\small\thepage}}
% la resta: el símbol a l'esquerra i el curs a la dreta, amb el títol al centre com abans
\fancyhead[L]{\raisebox{-0.3\height}{\includegraphics[height=0.62cm]{\fm@marca}}}
\fancyhead[R]{\small\bfseries\color{cursfosc}\fm@cursnom}
\renewcommand{\headrule}{\color{curs!45}\hrule height \headrulewidth}
\makeatother

\fmcurs{2eso}

\begin{document}
\thispagestyle{primera}

\capcalera{Matemàtiques}{2n ESO}{Fitxa d'exercicis · Unitat 1}{Nombres naturals i enters}
\vspace{10pt}
Nom i cognoms: \rule{8.4cm}{0.4pt}\hfill Grup: \rule{1.6cm}{0.4pt}\hfill Data: \rule{2.2cm}{0.4pt}
\vspace{8pt}

\begin{tcolorbox}[colback=soft,colframe=soft,arc=2mm,left=3mm,right=3mm,top=1.5mm,bottom=1.5mm]
\small Exercicis graduats per nivells de dificultat:\quad \estrelles{1}\ bàsic \qquad \estrelles{2}\ mitjà \qquad \estrelles{3}\ avançat.\\[2pt] Resol-los al quadern mostrant tots els passos. Fes els càlculs sense calculadora.
\end{tcolorbox}

% =====================================================
\seccio{1}{Els nombres naturals}

\ex{1} Calcula el quocient i el residu de cada divisió i fes la prova ($D=d\cdot q+r$):\begin{tasks}(4)
  \task $247:12$
  \task $1000:37$
  \task $505:5$
  \task $89:90$
\end{tasks}

\ex{1} Indica si el resultat és un nombre natural. Si no ho és, explica per què:\begin{tasks}(5)
  \task $7-12$
  \task $15\cdot 4$
  \task $20:6$
  \task $6:3+1$
  \task $3-3$
\end{tasks}

\ex{2} En una divisió el divisor és 15, el quocient 23 i el residu 9. Quin és el dividend? Podria ser 17 el residu d'aquesta divisió?

\ex{2} Calcula de cap aplicant la propietat distributiva:\begin{tasks}(4)
  \task $7\cdot 99$
  \task $12\cdot 15+12\cdot 5$
  \task $35\cdot 102$
  \task $8\cdot (10+5)$
\end{tasks}


% =====================================================
\seccio{2}{Múltiples i divisors}

\ex{1} Escriu els cinc primers múltiples no nuls de:\begin{tasks}(3)
  \task $7$
  \task $12$
  \task $15$
\end{tasks}

\ex{1} Escriu tots els divisors de:\begin{tasks}(5)
  \task $24$
  \task $36$
  \task $45$
  \task $60$
  \task $17$
\end{tasks}

\ex{2} Digues si són vertaderes o falses i justifica-ho:\begin{tasks}(1)
  \task 48 és múltiple de 6.
  \task 7 és divisor de 77.
  \task 5 és múltiple de 15.
  \task 1 és divisor de qualsevol nombre.
  \task 0 és múltiple de 9.
\end{tasks}


% =====================================================
\seccio{3}{Criteris de divisibilitat. Nombres primers}

\ex{2} Sense fer la divisió, indica per quins d'aquests nombres és divisible cadascun: 2, 3, 5, 9, 10 i 11.\begin{tasks}(4)
  \task $1320$
  \task $4455$
  \task $7125$
  \task $2002$
\end{tasks}\recorda{Per l'11: resta la suma de les xifres de lloc senar i la de les de lloc parell; ha de donar 0 o un múltiple d'11.}

\ex{2} Quines xifres poden anar en el lloc de $\square$?\begin{tasks}(1)
  \task $4\square7$ ha de ser múltiple de 3.
  \task $23\square$ ha de ser múltiple de 2 i de 5.
  \task $5\square2$ ha de ser múltiple de 9.
\end{tasks}

\ex{1} Classifica en primers i compostos. Dels compostos, escriu-ne la descomposició:\begin{tasks}(8)
  \task $31$
  \task $51$
  \task $57$
  \task $61$
  \task $87$
  \task $91$
  \task $97$
  \task $1$
\end{tasks}


% =====================================================
\seccio{4}{Descomposició en factors primers}

\ex{1} Descompon en factors primers:\begin{tasks}(6)
  \task $60$
  \task $84$
  \task $126$
  \task $360$
  \task $1000$
  \task $1176$
\end{tasks}

\ex{2} Quin nombre té cada descomposició?\begin{tasks}(3)
  \task $2^3\cdot3\cdot5^2$
  \task $2^2\cdot3^3$
  \task $3\cdot5\cdot7\cdot11$
\end{tasks}

\ex{3} Per quin nombre, el més xicotet possible, cal multiplicar 72 perquè el resultat siga un quadrat perfecte? De quin nombre és el quadrat?


% =====================================================
\seccio{5}{Mínim comú múltiple i màxim comú divisor}

\ex{1} Calcula llistant múltiples i divisors:\begin{tasks}(4)
  \task $\text{mcm}(12,18)$
  \task $\text{MCD}(12,18)$
  \task $\text{mcm}(8,20)$
  \task $\text{MCD}(8,20)$
\end{tasks}

\ex{2} Calcula el mcm i el MCD descomponent en factors primers:\begin{tasks}(3)
  \task $72$, $60$
  \task $84$, $126$
  \task $45$, $60$, $75$
  \task $140$, $300$
  \task $16$, $27$
\end{tasks}\recorda{\textbf{mcm}: factors comuns i no comuns amb el major exponent. \textbf{MCD}: només els comuns, amb el menor exponent.}

\ex{2} Comprova amb 72 i 60 que el producte del MCD pel mcm és igual al producte dels dos nombres.

\ex{3} El MCD de dos nombres és 6 i el mcm, 90. Si un dels nombres és 18, quin és l'altre?

\ex{2} Són primers entre si aquestes parelles de nombres?\begin{tasks}(3)
  \task $15$ i $28$
  \task $21$ i $35$
  \task $9$ i $16$
\end{tasks}


% =====================================================
\seccio{6}{Problemes de mcm i MCD}

\ex{2} Dues línies d'autobús ixen de la mateixa parada a les 7:00. La línia A hi passa cada 15 minuts i la B, cada 20 minuts. A quina hora tornaran a coincidir?

\ex{2} Una floristeria té 60 roses i 84 clavells. Vol fer el major nombre possible de rams iguals sense que sobre cap flor. Quants rams farà? Quantes roses i quants clavells durà cada ram?

\ex{2} Volem enrajolar el terra d'una habitació rectangular de 4,80 m × 3,60 m amb rajoles quadrades, les més grans possible i sense tallar-ne cap. Quant ha de mesurar el costat de cada rajola? Quantes rajoles calen?

\ex{2} Tres fars s'encenen cada 8, 12 i 18 segons, respectivament. Si ara s'acaben d'encendre alhora, quants segons passaran fins que tornen a encendre's junts?

\ex{2} Un magatzem té 150 llibres de matemàtiques i 210 de ciències. Es volen fer paquets iguals, amb el màxim nombre de llibres i amb un sol tipus de llibre en cada paquet. Quants llibres durà cada paquet? Quants paquets es faran?

\ex{3} Si agrupem els caramels d'una bossa de 6 en 6, de 8 en 8 o de 10 en 10, sempre en sobren 3. Quants caramels hi ha, com a mínim, a la bossa (més de 3)?

\ex{3} Dues rodes dentades engranades tenen 24 i 36 dents. Quantes voltes ha de fer cadascuna perquè tornen a estar en la posició inicial?


% =====================================================
\seccio{7}{Els nombres enters}

\ex{1} Expressa amb un nombre enter:\begin{tasks}(3)
  \task 3 graus baix zero
  \task el segon soterrani
  \task un deute de 50 €
  \task 1200 m d'altitud
  \task l'any 44 a. C.
  \task guanyar 15 €
\end{tasks}

\ex{1} Escriu el valor absolut i l'oposat de:\begin{tasks}(4)
  \task $-7$
  \task $12$
  \task $0$
  \task $-25$
\end{tasks}

\ex{1} Representa en la recta numèrica i ordena de menor a major:\[-3,\quad 5,\quad {-8},\quad 0,\quad 2,\quad {-1},\quad {-12},\quad 7\]

\ex{2} Escriu tots els nombres enters $x$ que compleixen:\begin{tasks}(4)
  \task $-4<x\le2$
  \task $\lvert x\rvert <3$
  \task $\lvert x\rvert =5$
  \task $-2\le x<-1$
\end{tasks}

\ex{2} Les temperatures mínimes d'una setmana han sigut: dilluns $-3$\,°C, dimarts $-5$\,°C, dimecres $1$\,°C, dijous $0$\,°C, divendres $-2$\,°C, dissabte $4$\,°C i diumenge $-1$\,°C. Quin dia va fer més fred? Ordena-les de menor a major. Quina diferència hi ha entre la més alta i la més baixa?


% =====================================================
\seccio{8}{Operacions amb enters}

\ex{1} Calcula:\begin{tasks}(4)
  \task $(-7)+(-5)$
  \task $(-9)+4$
  \task $12+(-15)$
  \task $(-6)-(-10)$
  \task $8-(+13)$
  \task $(-3)-5$
  \task $0-(-7)$
  \task $(-11)+11$
\end{tasks}

\ex{1} Calcula agrupant els positius i els negatius:\begin{tasks}(2)
  \task $5-8+3-7+2$
  \task $-4-6+10-3$
  \task $12-15-9+20$
  \task $-1+2-3+4-5+6$
\end{tasks}

\ex{2} Lleva els parèntesis i els claudàtors i calcula:\begin{tasks}(2)
  \task $7-(3-5+2)$
  \task $-(4-9)+(-2+6)$
  \task $10-\left[3-(5-8)\right]$
  \task $-\left[-(2-7)-(-3+1)\right]$
\end{tasks}\recorda{Un parèntesi precedit d'un signe $-$ es lleva canviant el signe de tots els termes de dins.}

\ex{1} Calcula aplicant la regla dels signes:\begin{tasks}(4)
  \task $(-6)\cdot (-4)$
  \task $7\cdot (-8)$
  \task $(-45):9$
  \task $(-56):(-7)$
  \task $0:(-5)$
  \task $(-2)\cdot (-3)\cdot (-5)$
  \task $(-1)\cdot (-1)\cdot (-1)\cdot (-1)$
  \task $72:(-8)$
\end{tasks}

\ex{2} Quin nombre falta en cada cas?\begin{tasks}(3)
  \task $(-4)\cdot\square=28$
  \task $\square:(-6)=5$
  \task $\square\cdot(-9)=-81$
  \task $(-48):\square=-12$
  \task $(-5)+\square=-12$
  \task $\square-8=-3$
\end{tasks}

\ex{2} Un submarí està a 120 m de profunditat. Puja 45 m i després baixa 80 m. A quina profunditat es troba?

\ex{2} Pitàgores va nàixer l'any 570 a. C. i va morir l'any 495 a. C. Quants anys va viure?

\ex{3} Un compte bancari té un saldo de 230 €. Es paguen tres rebuts de 85 € cadascun i, després, s'hi ingressen 120 €. Quin és el saldo final? Ha estat el compte en negatiu en algun moment?


% =====================================================
\seccio{9}{Operacions combinades}

\ex{1} Calcula respectant la jerarquia de les operacions:\begin{tasks}(3)
  \task $5-3\cdot 4$
  \task $(5-3)\cdot 4$
  \task $18:(-3)+2\cdot (-5)$
  \task $-20:4\cdot (-2)$
  \task $3\cdot (-2)-4\cdot (-3)$
\end{tasks}\recorda{Primer, parèntesis; després, potències i arrels; després, multiplicacions i divisions (d'esquerra a dreta) i, al final, sumes i restes.}

\ex{2} Calcula:\begin{tasks}(2)
  \task $4-2\cdot \left[3-(6-10)\right]$
  \task $(-3)\cdot (2-7)+15:(-5)$
  \task $\left[(-8)+5\right]\cdot \left[6-(-2)\right]$
  \task $-6-\left[4-3\cdot (2-5)\right]$
  \task $25:(-5)-(-3)\cdot (4-6):2$
\end{tasks}

\ex{3} Un alumne ha fet aquesta operació: \ $12-4:2\cdot(-3)=8:2\cdot(-3)=4\cdot(-3)=-12$. \ Explica quin error ha comés i calcula-la bé.

\ex{3} Posa parèntesis perquè les igualtats siguen certes:\begin{tasks}(3)
  \task $6-4\cdot3+2=10$
  \task $6-4\cdot3+2=-14$
  \task $-8:2-6=2$
\end{tasks}


% =====================================================
\seccio{10}{Potències i arrels amb enters}

\ex{1} Calcula:\begin{tasks}(4)
  \task $(-3)^{4}$
  \task $-3^{4}$
  \task $(-2)^{5}$
  \task $(-1)^{100}$
  \task $(-1)^{25}$
  \task $(-10)^{3}$
  \task $0^{7}$
  \task $(-7)^{0}$
\end{tasks}\recorda{Base negativa: exponent parell $\to$ resultat positiu; exponent senar $\to$ resultat negatiu. Compte: $-3^4=-(3^4)$, però $(-3)^4=(-3)\cdot(-3)\cdot(-3)\cdot(-3)$.}

\ex{2} Expressa com una sola potència:\begin{tasks}(3)
  \task $(-2)^3\cdot(-2)^5$
  \task $(-5)^9:(-5)^6$
  \task $\left[(-3)^2\right]^4$
  \task $2^5\cdot(-3)^5$
  \task $(-12)^4:3^4$
  \task $7^3\cdot7^4:7^5$
\end{tasks}

\ex{2} Calcula aplicant les propietats de les potències:\begin{tasks}(4)
  \task $(-2)^7:(-2)^4$
  \task $\left[(-1)^3\right]^5\cdot(-1)^2$
  \task $5^4\cdot2^4$
  \task $(-6)^3:(-3)^3$
\end{tasks}

\ex{1} Calcula les arrels quadrades (si existeixen):\begin{tasks}(6)
  \task $\sqrt{81}$
  \task $\sqrt{144}$
  \task $\sqrt{-25}$
  \task $\sqrt{0}$
  \task $\sqrt{400}$
  \task $\sqrt{-1}$
\end{tasks}

\ex{2} Calcula l'arrel entera i el residu i fes la comprovació:\begin{tasks}(4)
  \task $\sqrt{50}$
  \task $\sqrt{90}$
  \task $\sqrt{130}$
  \task $\sqrt{200}$
\end{tasks}

\ex{2} Calcula:\begin{tasks}(2)
  \task $(-2)^{3}+3^{2}$
  \task $-2^{4}+(-2)^{4}$
  \task $5-2\cdot \sqrt{36}+(-1)^{7}$
  \task $(-3)^{2}\cdot 2-\sqrt{49}:(-7)$
  \task $\left[(-2)^{2}-3^{2}\right]^{2}:(-5)$
  \task $\sqrt{25-9}-(-3)^{3}:9$
\end{tasks}

\ex{3} Calcula pas a pas, respectant la jerarquia de les operacions:\begin{tasks}(1)
  \task $\left[(-2)^{3}-(-3)^{2}\right]:(-17)+\sqrt{64}\cdot (-1)^{5}$
  \task $-3^{2}-\left[(-4)^{2}:(-2)^{3}-\sqrt{81}\right]\cdot (-2)$
  \task $\sqrt{5^{2}-3^{2}}\cdot (-2)^{3}-\left[(-6)^{2}:(-3)^{2}-2^{4}\right]^{2}:(-9)$
  \task $\left[2^{5}:((-2)^{2}-12)-(-1)^{20}\right]\cdot \sqrt{49}$
  \task $(-10)^{3}:(-5)^{2}+4\cdot \left[\sqrt{100}-(-3)^{3}:(-9)\right]^{2}$
\end{tasks}\recorda{Resol primer els parèntesis i claudàtors de dins cap a fora. Dins de cada un, primer les potències i les arrels, després les multiplicacions i divisions i, al final, les sumes i restes.}

\ex{2} Sense fer el càlcul, digues quin signe té el resultat:\begin{tasks}(5)
  \task $(-7)^{12}$
  \task $(-3)^{15}$
  \task $-5^{4}$
  \task $-(-2)^{6}$
  \task $(-1)^{2026}$
\end{tasks}

\ex{3} Una parcel·la quadrada té una superfície de 1225\,m². Quant mesura el costat? Quants metres de tanca calen per a envoltar-la?

\vfill
\begin{center}\color{muted}\small\faCheckCircle[regular]\ \ Ànim! Si et costa algun exercici, repassa els apunts interactius de la unitat.\end{center}

\end{document}
